%%=============================================================================
%% Samenvatting
%%=============================================================================

% TODO: De "abstract" of samenvatting is een kernachtige (~ 1 blz. voor een
% thesis) synthese van het document.
%
% Een goede abstract biedt een kernachtig antwoord op volgende vragen:
%
% 1. Waarover gaat de bachelorproef?
% 2. Waarom heb je er over geschreven?
% 3. Hoe heb je het onderzoek uitgevoerd?
% 4. Wat waren de resultaten? Wat blijkt uit je onderzoek?
% 5. Wat betekenen je resultaten? Wat is de relevantie voor het werkveld?
%
% Daarom bestaat een abstract uit volgende componenten:
%
% - inleiding + kaderen thema
% - probleemstelling
% - (centrale) onderzoeksvraag
% - onderzoeksdoelstelling
% - methodologie
% - resultaten (beperk tot de belangrijkste, relevant voor de onderzoeksvraag)
% - conclusies, aanbevelingen, beperkingen
%
% LET OP! Een samenvatting is GEEN voorwoord!

%%---------- Nederlandse samenvatting -----------------------------------------
%
% TODO: Als je je bachelorproef in het Engels schrijft, moet je eerst een
% Nederlandse samenvatting invoegen. Haal daarvoor onderstaande code uit
% commentaar.
% Wie zijn bachelorproef in het Nederlands schrijft, kan dit negeren, de inhoud
% wordt niet in het document ingevoegd.

% \IfLanguageName{english}{%
% \selectlanguage{dutch}
% \chapter*{Samenvatting}
% \lipsum[1-4]
% \selectlanguage{english}
% }{}

%%---------- Samenvatting -----------------------------------------------------
% De samenvatting in de hoofdtaal van het document

\chapter*{\IfLanguageName{dutch}{Samenvatting}{Abstract}}

De transparantie over het voedingsaanbod in Vlaamse supermarkten is beperkt, hoewel gezonde en duurzame voeding steeds centraler staan in beleid en onderzoek. 
Er bestaat geen geautomatiseerde manier om het werkelijke assortiment in kaart te brengen, waardoor onderzoekers, beleidsmakers en consumenten onvoldoende zicht hebben op de aanwezigheid van productcategorieën, duurzame alternatieven en marktverhoudingen. 
Dit onderzoek evalueert bestaande computer vision-modellen (bijv. YOLO-achtige object detectors, multimodale OCR- en beeld-classificatiemodellen) op de taak van automatische productherkenning in beeldmateriaal uit Vlaamse supermarkten, met als doel de nauwkeurigheid van detectie te beoordelen en eventueel te verbeteren via fine-tuning.
Het doel is een reproduceerbare methodologie te ontwikkelen die beelddata verwerkt en productinformatie herkent. 
De aanpak omvat het verzamelen van representatieve beelddata, het testen van modellen zoals YOLOv8 en Grounding DINO, het analyseren van prestaties via mAP en IoU-gebaseerde detectiemetrics, aangevuld met SKU-niveau Top-1 en Top-5 classificatie-accuratesse, macro-F1-scores en inventarisatiegerichte indicatoren zoals product coverage rate en assortimentsafwijking. 
Verwacht wordt dat deze studie de haalbaarheid aantoont van automatische inventarisatie in een complex retailecosysteem en dat de resultaten bijdragen aan betere monitoring van duurzaamheid en gezondheid, met concrete meerwaarde voor onderzoekers, beleidsmakers, consumenten en retailers.

